\documentclass[11pt]{article}

\usepackage[margin=1in]{geometry}
\usepackage{booktabs}
\usepackage{longtable}
\usepackage{array}
\usepackage{enumitem}
\usepackage{hyperref}

% simple, share-friendly formatting
\setlength{\parskip}{0.5em}
\setlength{\parindent}{0pt}
\renewcommand{\arraystretch}{1.2}

\title{Verdict Likert Dimensions\\[0.3em]
\large Current set, proposed extensions, and what the last run showed}
\author{}
\date{}

\begin{document}
\maketitle

Companion to \texttt{verdict\_confidence\_design.md} (the binary-verdict + confidence
redesign) and \texttt{confidence\_metrics.md} (the calibration metric suite). This document
is the single reference for \textbf{the score dimensions themselves}: what we score today
(verbatim from the prompt), how to grow to \textbf{3 truth + 3 confidence} dimensions, and
what the last run's distributions imply.

\paragraph{TL;DR}
\begin{itemize}[leftmargin=1.4em]
  \item Today the verifier scores \textbf{1 truth dimension} (\texttt{veracity}) +
    \textbf{3 confidence dimensions} (\texttt{evidence\_sufficiency},
    \texttt{evidence\_agreement}, \texttt{source\_reliability}), each 1--5.
  \item To reach \textbf{3 truth + 3 confidence}, add \textbf{2 truth dimensions}:
    \textbf{Contextual Integrity (misleadingness)} and \textbf{Attribution Fidelity} ---
    both orthogonal to \texttt{veracity}, both measurable text-only, both mapping to
    problems we already care about (the \S3b post-flag step and the attribution-safety rule).
  \item The 3 confidence dims are not equal: \texttt{evidence\_sufficiency} is the best
    evidence-side signal; \texttt{source\_reliability} carries \textbf{almost no signal}
    last run (AUROC 0.56; the literature agrees source-credibility is a weak heuristic)
    $\rightarrow$ \textbf{redefine toward evidence independence/provenance, or drop it.}
    And the model's own verdict-decisiveness is \textbf{anti-predictive} --- the usable
    confidence signal is \textbf{cross-model agreement} (see \S5).
\end{itemize}

\hrulefill

\section{The current dimensions (verbatim from \texttt{verify\_prompts.py::LIKERT\_SYSTEM})}

One LLM call scores the synthesized analysis. Prompt header: \emph{``Output one veracity
score (the verdict --- which way the evidence points) and three confidence scores (how much
to trust that read). Score each 1-5 independently against its own definition. Use the full
scale, including 2 and 4 deliberately.''} Parsed by \texttt{parse\_likert}, clamped 1--5,
stored in \texttt{models.py::VerdictScores}.

\subsection*{\texttt{veracity} --- truth (the verdict direction)}
\emph{Which way, and how strongly, does the evidence point on the claim's MAIN assertion?}

\begin{longtable}{@{}c>{\raggedright\arraybackslash}p{0.9\linewidth}@{}}
\toprule
5 & Clearly true --- evidence confirms the main assertion and its load-bearing specifics
(numbers, dates, named entities). \\
4 & Mostly true --- the main assertion holds, but a load-bearing specific may be approximate
or slightly off. \\
3 & Mixed or indeterminate --- neither clearly confirms nor clearly contradicts the main
assertion. \\
2 & Mostly false --- the main assertion is contradicted, though a detail may be unclear. \\
1 & Clearly false --- the evidence contradicts the main assertion. \\
\bottomrule
\end{longtable}

\subsection*{\texttt{evidence\_sufficiency} --- confidence}
\emph{How directly and fully does the evidence address the claim?}

\begin{longtable}{@{}c>{\raggedright\arraybackslash}p{0.9\linewidth}@{}}
\toprule
5 & Speaks directly to the main assertion AND its load-bearing specifics. \\
4 & Directly addresses the main assertion, but leaves a specific (a number, date, qualifier)
untouched. \\
3 & Partial --- touches the topic but addresses the specific assertion only loosely, or
through a single piece of evidence. \\
2 & Only topically adjacent --- does not speak to the specific assertion. \\
1 & Nothing in the analysis bears on the claim's substance. \\
\bottomrule
\end{longtable}

\subsection*{\texttt{evidence\_agreement} --- confidence}
\emph{Do the pieces of evidence point the same way?}

\begin{longtable}{@{}c>{\raggedright\arraybackslash}p{0.9\linewidth}@{}}
\toprule
5 & All pieces point the same way; no contradictions. \\
4 & Broad convergence; only minor differences of scope or wording. \\
3 & Some tension between pieces --- or only one piece exists, so agreement cannot be judged. \\
2 & Notable conflict --- at least one piece points each way. \\
1 & The evidence flatly contradicts itself on a load-bearing point. \\
\bottomrule
\end{longtable}

\subsection*{\texttt{source\_reliability} --- confidence}
\emph{How trustworthy are the cited sources, regardless of what they say?}

\begin{longtable}{@{}c>{\raggedright\arraybackslash}p{0.9\linewidth}@{}}
\toprule
5 & Primary or authoritative --- official records, peer-reviewed work, named domain experts,
court/legislative documents. \\
4 & Established mainstream outlets or institutional reporting. \\
3 & Mixed --- some reputable, but leaning on secondary or lower-tier reporting. \\
2 & Mostly blogs, opinion, or low-traffic sites. \\
1 & Only fringe, anonymous, or unidentifiable sources --- or none usable. \\
\bottomrule
\end{longtable}

\emph{(History: the earlier set was \texttt{veracity / coverage / consistency /
source\_quality}; coverage$\rightarrow$evidence\_sufficiency,
consistency$\rightarrow$evidence\_agreement, source\_quality$\rightarrow$source\_reliability.
The binary verdict (\texttt{pass} iff supported, else \texttt{flag}) is derived from
\texttt{veracity} by a threshold --- see \texttt{verdict\_confidence\_design.md}.)}

Two related prompts score the same analysis: a 4-class evaluator (ClaimCheck verbatim:
Supported / Refuted / Conflicting Evidence-Cherrypicking / Not Enough Evidence) and the
\S3b post-level FLAG step (judges the raw post for misleadingness given the atomic verdict).

\hrulefill

\section{Why decompose ``truth'' into more than one axis}

\texttt{veracity} is a single true$\leftrightarrow$false spine. Every mature fact-checking
scheme refuses to live on that spine alone --- they either fold context into the mid-scale
or add \textbf{non-degree tags}:

\begin{longtable}{@{}>{\raggedright\arraybackslash}p{0.18\linewidth}
  >{\raggedright\arraybackslash}p{0.13\linewidth}
  >{\raggedright\arraybackslash}p{0.27\linewidth}
  >{\raggedright\arraybackslash}p{0.22\linewidth}
  >{\raggedright\arraybackslash}p{0.12\linewidth}@{}}
\toprule
\textbf{Scheme} & \textbf{Truth-degree} & \textbf{Context/framing} & \textbf{Attribution} &
\textbf{Temporal} \\
\midrule
PolitiFact (6-pt) & \checkmark & \textbf{\checkmark\ baked in} (Half True = ``out of
context'') & --- & implicit \\
WaPo Pinocchios & partial & \textbf{\checkmark} (2 = ``misleading impression\ldots playing
with words'') & --- & repetition axis \\
\textbf{Snopes} & \checkmark & tag: \emph{Miscaptioned} & \textbf{tags: \emph{Misattributed}
/ \emph{Correct Attribution}} & tag: \emph{Outdated} \\
schema.org ClaimReview & numeric + textual & free-text & --- & --- \\
FEVER & 3-way (+NEI) & --- & --- & --- \\
\textbf{AVeriTeC} & 4-way & \textbf{\checkmark} Cherrypicking is a verdict class & --- & --- \\
Wardle 7 types & --- & Misleading / False-context / False-connection & Imposter /
misattribution & False context \\
\bottomrule
\end{longtable}

Journalistic scales (PolitiFact/WaPo/Snopes) separate literal truth from context,
attribution, and time; Snopes makes attribution and time \textbf{first-class tags orthogonal
to its truth scale} --- the cleanest existence proof that these are distinct axes, not
sub-cases of degree-of-truth. A data-driven mapping study finds these scales only align once
collapsed to binary real/fake, and that ``Misleading/Unproven/Miscaptioned'' categories
don't map onto a linear truth line. That is the empirical case for \textbf{adding axes rather
than stretching \texttt{veracity}.}

\hrulefill

\section{Proposed 2 new truth dimensions}

\subsection*{Truth dim \#2 --- Contextual Integrity (misleadingness / framing fidelity)}
\emph{Beyond the literal assertion, does the impression the post creates survive the full
evidence --- or do omitted context, cherry-picking, or implied causation make a
literally-true claim mislead?}

\begin{longtable}{@{}c>{\raggedright\arraybackslash}p{0.9\linewidth}@{}}
\toprule
5 & Literal content \textbf{and} its evident implication are supported; no material context
omitted that would change a reader's takeaway. \\
3 & Literally accurate but missing context/qualification a reasonable reader needs; a
\emph{partially} unsupported impression ($\approx$ PolitiFact \textbf{Half True}). \\
1 & True-ish details arranged to imply a conclusion the evidence contradicts ---
cherry-picking, implied causation, \textbf{stale event shown as current} ($\approx$ paltering
/ \textbf{Cherrypicking}). \\
\bottomrule
\end{longtable}

\begin{itemize}[leftmargin=1.4em]
  \item \textbf{Why:} this is the single largest gap between \texttt{veracity} and ground
    truth, and it is exactly what the \S3b post-flag step needs to judge. Fact-checkers treat
    it as core (PolitiFact's busiest mid-ratings are \emph{defined} by context omission;
    AVeriTeC promotes Cherrypicking to a verdict class). The psychology is named:
    \textbf{paltering} = ``deceiving through truthful statements.''
  \item \textbf{Distinct from} \texttt{veracity} (scores the literal assertion; CI scores
    whether the \emph{gist} holds even when the literal claim is true) and from
    \texttt{evidence\_*} (those are about the evidence we gathered; CI is about
    claim-framing-vs-evidence).
  \item \textbf{Measurable text-only:} yes --- we already retrieve the surrounding evidence.
    The added step: (i) state the literal proposition, (ii) name the \emph{implied}
    proposition, (iii) test whether full evidence supports the implication / whether a
    material fact was omitted. \textbf{Anchor it to hold literal truth fixed} (``literally
    accurate \textbf{but}\ldots'') so it can't silently re-score \texttt{veracity}.
  \item \textbf{Folds in temporal:} the ``stale-as-current'' case is encoded as the level-1
    anchor for now (see \S5 on Temporal Validity as a later standalone axis).
\end{itemize}

\subsection*{Truth dim \#3 --- Attribution Fidelity}
\emph{For ``X said/did Y'' claims, does the evidence confirm X is correctly the source/agent
of Y --- scored independently of whether Y itself is true.}

\begin{longtable}{@{}c>{\raggedright\arraybackslash}p{0.9\linewidth}@{}}
\toprule
5 & Evidence confirms the named source/actor said/did exactly what's claimed, in context. \\
3 & Roughly the right source but the quote is altered, partial, or stripped of qualifying
context; or right words, wrong stance. \\
1 & Misattributed/fabricated quote, imposter source, or words so out of context the meaning
inverts ($\approx$ Snopes \textbf{Misattributed}; Wardle \textbf{imposter}). \\
N/A & Claim makes no attribution $\rightarrow$ \textbf{masked, not scored low} (see risks). \\
\bottomrule
\end{longtable}

\begin{itemize}[leftmargin=1.4em]
  \item \textbf{Why:} this directly implements our hard nudge-safety constraint --- \emph{do
    not flag accurate reporting of a false claim as if the poster lied.} ``X said Y'' is two
    checks (did X say it; is Y true). Snopes carries distinct \emph{Misattributed} /
    \emph{Correct Attribution} tags orthogonal to its truth scale; claim-normalization
    research explicitly keeps original quotes and decontextualizes the reference separately
    from verifying the content.
  \item \textbf{Distinct from} \texttt{veracity}: a post can be attribution 5 / veracity 1
    (faithfully quotes a false statement $\rightarrow$ should \textbf{pass} as accurate
    reporting). Not a statement about evidence quality.
  \item \textbf{Measurable text-only:} yes for quoted/attributed claims --- retrieval can
    locate the quote and check wording + context. \textbf{Needs an explicit ``applicable?''
    gate}; most claims are non-attributed and must be \emph{masked} (N/A), not penalized.
\end{itemize}

\hrulefill

\section{Re-assessing the 3 confidence dimensions}

\begin{itemize}[leftmargin=1.4em]
  \item \textbf{\texttt{evidence\_sufficiency} --- keep.} The most-attested confidence facet
    in claim verification (FEVER's \emph{NEI}, AVeriTeC's \emph{Not Enough Evidence} are
    dedicated classes), and it helped in our data. Our cleanest confidence signal.
  \item \textbf{\texttt{evidence\_agreement} --- keep, watch overlap.} Conflicting evidence
    is a first-class state. Risk: AVeriTeC couples agreement with \emph{cherry-picking},
    which is Contextual Integrity above --- keep them disjoint (\textbf{agreement =
    evidence-vs-evidence; CI = claim-framing-vs-evidence}).
  \item \textbf{\texttt{source\_reliability} --- redefine or replace.} It carried
    \textbf{almost no signal} last run (AUROC 0.56), consistent with the literature:
    source-credibility is a weak heuristic and predicting correctness from source features is
    empirically poor (esp. on the unreliable end). It is plausibly \textbf{confounded} ---
    hard/contested claims attract lower-grade sources, so low reliability co-occurs with
    hard-but-correct reads. \textbf{Options:} (a) redefine toward \textbf{evidence
    provenance/independence} (count of \emph{independent} corroborating sources,
    primary-vs-secondary) rather than reputational trust; or (b) drop it. The slot is better
    spent on \textbf{cross-model agreement} --- the single best confidence signal we measured
    (\S5), i.e.\ a self-consistency / panel facet, the best-attested confidence family.
\end{itemize}

\hrulefill

\section{Distribution \& calibration analysis (last run)}

\textbf{The collaborators' question --- did lower confidence track lower accuracy? Short
answer: no, not reliably, and the model's own decisiveness is the \emph{worst} confidence
signal.} Verified by re-running
\texttt{eval/scripts/verdict\_confidence/dimension\_distributions.py} (full tables in
\texttt{results/dimension\_analysis.md}). Data: \textbf{\texttt{agreement.parquet}} (n=60,
all 4 dims across a 4-model panel; \texttt{correct = deployment\_binary==gold}, acc 0.80) for
the multi-dim read; \textbf{\texttt{flag\_eval.parquet}} (n=1024, \texttt{veracity} +
\texttt{bare\_correct}, acc 0.85) for power. \emph{n=60 $\rightarrow$ panel numbers are
directional; the n=1024 results carry the powered claims.}

\paragraph{(a) The scale is barely used --- opposite failures on the two halves.}
The prompt says ``use 2 and 4 deliberately''; it largely doesn't. \texttt{veracity} collapses
to the \textbf{poles} (level-3 = 2.7\%, \textbf{81.8\% of mass at \{1,5\}}); the 3 confidence
dims pin to the \textbf{ceiling} (means 4.09--4.37, $\sim$85\% at $\geq$4). Little dynamic
range --- this is the \textbf{bigger calibration limiter} than the choice of dims.

\paragraph{(b) \texttt{veracity} is a direction, not a confidence --- and ``confident'' is
the high-risk bucket.} Errors don't sit at the uncertain middle; they cluster at
\textbf{confident-TRUE}:

\begin{longtable}{@{}>{\raggedright\arraybackslash}p{0.28\linewidth}
  >{\raggedright\arraybackslash}p{0.33\linewidth}
  >{\raggedright\arraybackslash}p{0.33\linewidth}@{}}
\toprule
 & \textbf{confident-FALSE (\texttt{veracity=1})} & \textbf{confident-TRUE
(\texttt{veracity=5})} \\
\midrule
flag\_eval (n=1024) & err \textbf{2.2\%} (n=544) & err \textbf{30.6\%} (n=294) --- \emph{all
90 errors are misinformation that PASSED} \\
panel (n=60) & err 7.1\% & err 54.5\% --- \emph{all errors gold=flag} \\
\bottomrule
\end{longtable}

Accuracy by \texttt{veracity} level (flag\_eval): v1 \textbf{0.978}, v2 0.719, v3 0.821, v4
0.652, v5 \textbf{0.694}; raw \texttt{veracity} correlates \textbf{$-0.36$} with correctness
--- a high ``true'' call is \emph{less} often right. The deployment headline: \textbf{never
trust a confident PASS at face value.}

\paragraph{(c) Which signals actually predict correctness} (panel n=60, AUROC / AURC;
base-err 0.200):

\begin{longtable}{@{}>{\raggedright\arraybackslash}p{0.42\linewidth}cc
  >{\raggedright\arraybackslash}p{0.32\linewidth}@{}}
\toprule
\textbf{Signal} & \textbf{AUROC} & \textbf{AURC} & \\
\midrule
cross-model agreement ($-$veracity std) & \textbf{0.658} & \textbf{0.113} & best --- monotone
(acc 0.65$\rightarrow$0.85$\rightarrow$0.90 by tertile) \\
\texttt{evidence\_sufficiency} mean & \textbf{0.650} & 0.117 & best evidence-side dim \\
panel veracity magnitude & 0.622 & 0.139 & \\
weakest-link \texttt{min} of conf dims & 0.616 & 0.128 & \\
\texttt{evidence\_agreement} mean & 0.598 & 0.152 & \\
\texttt{source\_reliability} mean & 0.560 & 0.153 & \textbf{near-useless} \\
ensemble verdict decisiveness $|$veracity$-3|$ & \textbf{0.486} & 0.218 &
\textbf{anti-predictive} (acc \emph{falls} 0.85$\rightarrow$0.80$\rightarrow$0.75) \\
\bottomrule
\end{longtable}

The only useful confidence signals are \textbf{cross-model agreement} (ask several models;
trust what they agree on) and \textbf{\texttt{evidence\_sufficiency}}. The verdict's own
decisiveness is worse than random ordering. \textbf{This does NOT replicate the earlier
``veracity-magnitude AUROC $\approx$0.76'' note} --- that was a single-model/dev-slice read;
on the reproducible run data magnitude is a poor signal. Flag for re-confirmation on a larger
panel.

\paragraph{(d) The dims are NOT redundant.} Confidence-dim Spearman $\rho$ = 0.33--0.59 (all
well below 0.9) $\rightarrow$ the decomposition is real; keep all three. Each confidence dim
is \emph{negatively} correlated with \texttt{veracity} ($-0.10$ to $-0.40$) --- the model
rates evidence ``sufficient/agreeing'' more when it leans false.

\paragraph{Implications:} (i) fix the scale collapse first (few-shot anchors / force 2--4) ---
it caps everything downstream; (ii) gate deployment confidence on \textbf{agreement +
evidence\_sufficiency}, not the verdict's decisiveness; (iii) treat any confident PASS as the
high-risk bucket --- exactly what the new \textbf{Contextual Integrity} + \textbf{Attribution
Fidelity} truth dims are meant to catch (a confident-true literal claim that is misleading or
misattributed).

\hrulefill

\section{Proposed final dimension set (3 truth + 3 confidence)}

\begin{longtable}{@{}>{\raggedright\arraybackslash}p{0.13\linewidth}
  >{\raggedright\arraybackslash}p{0.32\linewidth}
  >{\raggedright\arraybackslash}p{0.48\linewidth}@{}}
\toprule
\textbf{Class} & \textbf{Dimension} & \textbf{Status} \\
\midrule
Truth & \texttt{veracity} & keep (drives the binary verdict) \\
Truth & \textbf{Contextual Integrity} & \textbf{new} --- misleadingness/framing; powers the
\S3b flag \\
Truth & \textbf{Attribution Fidelity} & \textbf{new} --- masked when no attribution; powers
nudge-safety \\
Confidence & \texttt{evidence\_sufficiency} & keep (strongest confidence facet) \\
Confidence & \texttt{evidence\_agreement} & keep (keep disjoint from Contextual Integrity) \\
Confidence & \texttt{source\_reliability} $\rightarrow$ \textbf{provenance/independence}
\emph{or} \textbf{self-consistency} & redefine or replace (anti-predictive as-is) \\
\bottomrule
\end{longtable}

\textbf{Gate every new axis on a validation check} before shipping: does it add
\emph{marginal} AUROC over the existing set, and does its learned weight have the expected
sign? \texttt{source\_reliability} is the cautionary tale --- a plausible axis that encoded
claim-hardness, not trustworthiness.

\hrulefill

\section{Open decisions (for the meeting)}

\begin{itemize}[leftmargin=1.4em]
  \item The two new truth dims --- agree on \textbf{Contextual Integrity + Attribution
    Fidelity}, or swap in \textbf{Temporal Validity} (real but harder to measure text-only;
    recommended as a phase-2 promotion from CI's stale-as-current anchor).
  \item Does \texttt{veracity} stay one 1--5 verdict score, or do the 3 truth dims
    \textbf{compose} into the verdict?
  \item \texttt{source\_reliability}: redefine (provenance/independence) vs replace
    (self-consistency) vs drop.
  \item Few-shot vs the current zero-shot Likert prompt (the ``use 2 and 4'' instruction is
    currently zero-shot --- the scale-usage numbers in \S5 will say whether it's obeyed).
\end{itemize}

\hrulefill

\section*{References}

\begin{itemize}[leftmargin=1.4em]
  \item PolitiFact methodology (context baked into mid-ratings) $\cdot$ WaPo Pinocchios
    $\cdot$ Snopes ratings (Misattributed/Miscaptioned/Outdated tags).
  \item Wardle, \emph{Fake news. It's complicated.} (First Draft 7 types).
  \item Schlichtkrull et al., \textbf{AVeriTeC} (FEVER 2024) --- Cherrypicking \& NEI as
    first-class verdicts.
  \item Rogers et al., \textbf{paltering} (deception via truthful statements), PMC7259623.
  \item \textbf{ChronoFact} / Evidence-Based Temporal Fact Verification (arXiv:2407.15291)
    --- temporal axis.
  \item Source-reliability prediction is weak (arXiv:2410.18803);
    source-credibility-as-heuristic (Liu et al.\ 2025, \emph{Communication Research}).
  \item Claim Normalization (EMNLP'23 Findings); Document-level Claim Extraction \&
    Decontextualisation (arXiv:2406.03239).
\end{itemize}

\emph{Full annotated survey (landscape table, rankings, all citations): the research brief
feeding this doc; calibration metric definitions in \texttt{confidence\_metrics.md}.}

\end{document}
